%! @@TITLE@@ — Unit @@UNITINT@@, Lesson @@LESSONID@@ — Guided Notes & Practice (Answer Key)
% Copy the blank component VERBATIM, then: swap -boxes for -key, retitle the pageheader,
% replace each \blank{} with a short \ans{}, each \writelines{n} with exactly n \ansline{}
% calls, and each \termblank{T} with \termans{T}{definition} (TWO arguments, defined in
% @@PREFIX@@-key.sty — never redefine it here; it reuses \termblank's fixed row height, so a
% definition that wraps to a second line still costs exactly what the blank reserved).
% Every \begin{work} block and every \boxguard stays BYTE-IDENTICAL to the blank's.
% Keep \ans{} texts short enough not to wrap wider than the blank they replace, and keep
% each \ansline{} short enough not to wrap past the line it fills — that is what keeps the
% key page-for-page with the blank.
% \ans{} IS TEXT MODE: never inside $...$, and wrap \sqrt/\dfrac/\hat/\bar in $...$ inside it.
% NO teachernote here — teacher prose goes in the lesson plan.
\documentclass[12pt]{article}
\usepackage{@@PREFIX@@-article}
\usepackage{@@PREFIX@@-key}   % pulls in -boxes; do NOT also load -boxes

\begin{document}

\pageheader{Unit @@UNITINT@@, Lesson @@LESSONID@@}{Guided Notes \& Practice --- Answer Key}
% NAMESTRIP: no \namedateperiod here — mirrors the blank.

\vspace{0.12in}

% TODO: mirror the blank component exactly. Examples:
%   \termblank{Variable}  ->  \termans{Variable}{a characteristic recorded for every individual; a column heading, never an entry under it.}
%   \blank{2.4cm}         ->  \ans{context}
%   \writelines{2}        ->  \ansline{first line of the answer}
%                             \ansline{second line of the answer}

\end{document}
